% Source PDF: source/普通高考/2012/2012湖北理.pdf, page 1, question 12.
% PDF embedded bitmap attempt: pdfimages -list returned no image objects; the flowchart is vector page content.
% Grid evidence: tmp/review_2012_hubei_science/grid/page1_grid100.jpg and q12_grid.jpg.
% Original-side comparison crop: tmp/review_2012_hubei_science/crops/q12_original.png from a no-grid page render.
% Elements: rounded start/end terminals, initialization and update process boxes,
% n<3 decision diamond, output parallelogram, directed connectors and loop-back branch.
% Node-edge list: start->(a=1,s=0,n=1)->(s=s+a)->(a=a+2)->(n<3);
% (n<3)-是->(n=n+1)->loop-back to the stem before (s=s+a);
% (n<3)-否->输出s->结束. No other edges or dashed segments.
% Node text is centered with local zero paragraph indent; the loop return arrow
% merges into the main stem before the summation box and never enters output.
\begin{tikzpicture}[
  >=Stealth,
  x=1cm,y=1cm,
  line width=0.45pt,
  line cap=round,line join=round,
  font=\normalsize,
  flowtext/.style={anchor=center,align=center,inner sep=0pt,
    execute at begin node={\setlength{\parindent}{0pt}}},
  flowbox/.style={flowtext,draw,minimum height=.68cm,inner xsep=4pt,inner ysep=2pt},
  terminal/.style={flowbox,rounded corners=.16cm,minimum width=1.25cm},
  process/.style={flowbox},
  io/.style={flowbox,shape=trapezium,trapezium left angle=72,trapezium right angle=108,
    minimum height=.74cm},
  decision/.style={flowbox,diamond,aspect=2.8,minimum width=3.15cm,minimum height=.88cm,
    inner sep=0pt}
]
  \path[use as bounding box] (-2.85,-7.30) rectangle (4.45,0.42);

  \node[terminal] (start) at (0,0) {开始};
  \node[process,minimum width=3.45cm] (init) at (0,-1.12) {$a=1,s=0,n=1$};
  \node[process,minimum width=1.72cm] (sum) at (0,-2.28) {$s=s+a$};
  \node[process,minimum width=1.72cm] (incA) at (0,-3.38) {$a=a+2$};
  \node[decision] (test) at (0,-4.62) {$n<3?$};
  \node[process,minimum width=1.82cm] (incN) at (0,-5.80) {$n=n+1$};
  \node[io,minimum width=1.85cm] (output) at (2.80,-5.02) {输出 $s$};
  \node[terminal] (finish) at (2.80,-6.20) {结束};

  % Main chain and false/output branch.
  \draw[->] (start.south) -- (init.north);
  \draw[->] (sum.south) -- (incA.north);
  \draw[->] (incA.south) -- (test.north);
  \draw[->] (test.east) -- node[flowtext,above] {否} (2.80,-4.62) -- (output.north);
  \draw[->] (output.south) -- (finish.north);

  % True branch updates n and returns left/up to the stem before s=s+a.
  \draw[->] (test.south) -- node[flowtext,right] {是} (incN.north);
  \coordinate (loop-left) at (-2.35,-6.14);
  \coordinate (merge-level) at (0,-1.55);
  \draw (incN.south) -- (loop-left) -- (loop-left |- merge-level);
  \draw[->] (loop-left |- merge-level) -- (merge-level);
  \draw (init.south) -- (merge-level);
  \draw[->] (merge-level) -- (sum.north);

\end{tikzpicture}
